\input{preamble}

\begin{document}

\title{Deformation theory}
\maketitle

\phantomsection
\label{section-phantom}

\tableofcontents

\section{Introduction}
\label{section-intro}

\noindent
\cite{Sernesi-Overview} has a great way of
motivating deformation
theory via the moduli problem.
Here let's just put the following important
meaning for
``deformation theory'' (as opposed to
studying flat families!):

\begin{quote}
``{\it Deformation theory} is the study of
infinitesimal deformaitons
as a tool to understand the local structure
of the moduli space.
The goal is to be able to describe the
restriction
of the universal family to a small
neighborhood of $m \in \mathcal{M}$,
or, more precisely, its restriction to the
germ of $M$ at $m$.

What is interesting here is that we can study
order and infinitesimal
deformations even though the functor $F$ is
not representable or simply we don’t
yet know it is. This is the most frequent
case. Such an investigation will
reveal the infinitesimal properties at $[X]$
of a yet unknown global structure
on $M$ which will be hopefully understood at
a subsequent stage of the
investigation.  In order words it turns out
to be possible and convenient to
separate the {\it global moduli problem} from
the {\it local moduli problem},
and deformation theory studies the latter,
with the purpose of constructing a
family of deformations of a given object
parametrized by the spectrum of a local
ring, and having properties as close as
possible to a universal property.''
\end{quote}


\section{A short history of deformation theory}
\label{section-history-deformation-theory}

\noindent
The following chronology explains where the
Hilbert scheme, deformation functors and
miniversal deformation spaces enter the
story.

\begin{enumerate}
\item From 1949 to 1954 Grothendieck worked
mainly in functional analysis: topological
tensor products, nuclear spaces and the
theory of Banach spaces. By the end of 1954
he had decided to leave this subject. In
1955 his interests shifted first to
homological algebra and then increasingly to
algebraic geometry. See the
\href{https://research-portal.uu.nl/ws/portalfiles/portal/254168358/Life_and_Work_of_Alexander_Grothendieck.pdf}
{biographical account of Chai and Oort}.

\item In 1957 he obtained the
Grothendieck--Riemann--Roch theorem. At the
1958 International Congress of
Mathematicians he announced his program for
the cohomology of abstract algebraic
varieties. Thus his dream of a general
cohomological framework was already present
before the construction of the Hilbert
scheme. See
\href{https://webusers.imj-prg.fr/~leila.schneps/grothendieckcircle/CohomologyVarieties.pdf}
{Grothendieck's 1958 ICM paper}.

\item In 1960 the first volume of EGA,
\emph{Le langage des sch\'emas}, appeared.
Algebraic geometry was of course already an
old subject; what Grothendieck was creating
was its modern scheme-theoretic language.

\item In May 1961 Grothendieck proved the
representability of the Hilbert functor. More
precisely, if $X\to S$ is projective and a
Hilbert polynomial $P$ is fixed, then the
functor of flat families of closed
subschemes of $X$ with Hilbert polynomial
$P$ is represented by a projective
$S$-scheme $\operatorname{Hilb}^P_{X/S}$.
See
\href{https://www.numdam.org/item/SB_1960-1961__6__249_0/}
{Grothendieck's Bourbaki expos\'e 221}.

This result belonged to his wider program of
\emph{construction theorems}: formulate a
geometric classification problem as a
functor, and then prove that the functor is
represented by a scheme. Hilbert schemes,
Quot schemes and Picard schemes are central
examples of this philosophy.

\item There was simultaneously an analytic
line of development. Kodaira and Spencer
developed the deformation theory of compact
complex manifolds in 1958, and in 1962
Kuranishi constructed locally complete, or
miniversal, analytic families.

\item In 1968 Schlessinger's
\emph{Functors of Artin Rings} gave precise
criteria for an infinitesimal deformation
functor to be prorepresentable, or at least
to possess a hull. The hull is precisely the
formal miniversal deformation space that we
shall define below. See
\href{https://doi.org/10.1090/S0002-9947-1968-0217093-3}
{Schlessinger's original paper}.

\item In 1969 Michael Artin proved his
approximation and algebraization theorems.
These results form a bridge from formal
moduli problems to actual algebraic
families. See
\href{https://www.numdam.org/articles/10.1007/BF02684596/}
{Artin's 1969 approximation paper}.
\end{enumerate}

\noindent
After 1961 Grothendieck continued developing
the EGA and SGA machinery: formal geometry,
topoi, descent, the \`etale fundamental
group and \`etale cohomology. The last of
these supplied the main cohomological
machinery for the Weil conjectures. His more
ambitious vision was the theory of
\emph{motives}: a universal geometric source
from which the different reasonable
cohomology theories should arise. This
program was not completed. In particular,
Grothendieck's standard conjectures on
algebraic cycles remain open. He left the
IH\'ES in 1970, and Deligne proved the final
Weil conjecture in 1974 by a different
route.

\begin{remark}[Algebraization in the
Gr\"unbaum paper]
Two related results must not be conflated.
Grothendieck's existence theorem, developed
in EGA III in 1961, gives effectivity and
algebraization results over a complete local
base in the proper or projective setting.
Artin's later approximation and
algebraization theorems allow one, under
suitable hypotheses, to pass from a formal
versal deformation to an algebraic family
of finite type. See
\href{https://www.numdam.org/item/?id=PMIHES_1961__11__5_0}
{EGA III} and the
\href{https://stacks.math.columbia.edu/tag/04UW}
{Stacks Project's historical summary}.

Consequently, the sentence in our current
handout saying that properness and
``Grothendieck's Existence Theorem'' turn a
formal family directly into a family over
an algebraic curve is probably too quick.
The passage from a formal family over
$\operatorname{Spf}k[[s]]$ to a family over
an algebraic curve generally involves Artin
approximation or algebraization. This step
must be stated carefully in the paper.
\end{remark}


\section{Structures over the Dual Numbers}
\label{section-structures-dual-numbers}

\noindent
Other than the contents of
\cite[Chapter 1, Section 2]
{Hartshorne-deformation}, which
is most of the content of this section,
let us recall from Ariel's
lecture that he proved a proposition
with 4 equivalent conditions
for flatness over the dual numbers:
being a free module, having
$\operatorname{Tor}^1=0$, and two more
involving annihilators.

\begin{definition}
\label{definition-dual-numbers-first-order}
The very first deformation question to study
is structures over the \emph{dual numbers}
$$
D=k[t]/t^2.
$$
That is, one gives a structure (e.g., a
scheme, or a scheme with a subscheme, or a
scheme with a sheaf on it) over $k$ and one
seeks to classify extensions of this
structure over the dual numbers. These are
also called \emph{first-order deformations}.
To ensure that our structure is evenly
spread out over the base, we will always
assume that the extended structure is flat
over $D$. Flatness is the technical condition
that corresponds to the intuitive idea of a
deformation.
\end{definition}

\begin{lemma}
\label{lemma-flatness-tor}
A module $M$ over a noetherian ring $A$ is
flat if and only if for every prime ideal
$\mathfrak{p}\subseteq A$,
$$
\operatorname{Tor}^A_1
(M,A/\mathfrak{p})=0.
$$
\end{lemma}

\begin{proposition}
\label{proposition-square-zero-flatness}
Let $A'\to A$ be a surjective homomorphism
of noetherian rings whose kernel $J$ has
square zero. Then an $A'$-module $M'$ is flat
over $A'$ if and only if
\begin{enumerate}
\item
$$
M=M'\otimes_{A'}A
$$
is flat over $A$, and
\item the natural map
$$
M\otimes_AJ\longrightarrow M'
$$
is injective.
\end{enumerate}
\end{proposition}

\begin{definition}
\label{definition-embedded-first-order}
Let $X$ be a scheme over $k$ and let $Y$ be
a closed subscheme of $X$. We define a
\emph{deformation of $Y$ over $D$ in $X$} to
be a closed subscheme
$$
Y'\subseteq X'=X\times D,
$$
flat over $D$, such that
$$
Y'\times_D k=Y.
$$
\end{definition}

We consider the affine case first. Then $X$
corresponds to a $k$-algebra $B$, and $Y$ is
defined by an ideal $I\subseteq B$. We are
seeking ideals
$$
I'\subseteq B'=B[t]/t^2
$$
with $B'/I'$ flat over $D$ and such that the
image of $I'$ in $B=B'/tB'$ is just $I$.
Note that
$$
(B'/I')\otimes_Dk=B/I.
$$
Since $B$ is automatically flat over $k$,
by (2.2) the flatness of
$B'/I'$ over $D$ is equivalent to the
exactness of the sequence
$$
0\longrightarrow B/I
\xrightarrow{t}B'/I'
\longrightarrow B/I
\longrightarrow 0.
$$

\begin{proposition}
\label{proposition-affine-first-order}
In the situation above, to give
$I'\subseteq B'$ such that $B'/I'$ is flat
over $D$ and the image of $I'$ in $B$ is $I$
is equivalent to giving an element
$$
\varphi\in\operatorname{Hom}_B(I,B/I).
$$
In particular, $\varphi=0$ corresponds to the
trivial deformation given by
$$
I'=I\oplus tI
$$
inside $B'\cong B\oplus tB$.
\end{proposition}

\begin{definition}
\label{definition-normal-sheaf}
Before stating the conclusion, we will
define the \emph{normal sheaf of $Y$ in
$X$}. Note that the group
$$
\operatorname{Hom}_X
(\mathcal{I},\mathcal{O}_Y)
$$
can be regarded as
$$
H^0\bigl(X,
\mathcal{H}om_X
(\mathcal{I},\mathcal{O}_Y)\bigr).
$$
Furthermore, homomorphisms of $\mathcal{I}$
to $\mathcal{O}_Y$ factor through
$\mathcal{I}/\mathcal{I}^2$, which is a sheaf
on $Y$. So
$$
\mathcal{H}om_X
(\mathcal{I},\mathcal{O}_Y)
=
\mathcal{H}om_Y
(\mathcal{I}/\mathcal{I}^2,
\mathcal{O}_Y),
$$
and this latter sheaf is called the normal
sheaf of $Y$ in $X$, and is denoted by
$\mathcal{N}_{Y/X}$. If $X$ is nonsingular
and $Y$ is a locally complete intersection
in $X$, then $\mathcal{I}/\mathcal{I}^2$ is
locally free, so $\mathcal{N}_{Y/X}$ is
locally free also and can be called the
\emph{normal bundle of $Y$ in $X$}. This
terminology derives from the fact that if
$Y$ is also nonsingular, there is an exact
sequence
$$
0\longrightarrow\mathcal{T}_Y
\longrightarrow\mathcal{T}_X|_Y
\longrightarrow\mathcal{N}_{Y/X}
\longrightarrow 0,
$$
where $\mathcal{T}_Y$ and $\mathcal{T}_X$
denote the tangent sheaves to $Y$ and $X$,
respectively. In this case, therefore,
$\mathcal{N}_{Y/X}$ is the usual normal
bundle.
\end{definition}

\begin{theorem}
\label{theorem-embedded-first-order}
Let $X$ be a scheme over a field $k$, and let
$Y$ be a closed subscheme of $X$. Then the
deformations of $Y$ over $D$ in $X$ are in
natural one-to-one correspondence with
elements of $H^0(Y,\mathcal{N}_{Y/X})$, the
zero element corresponding to the trivial
deformation.
\end{theorem}

\begin{lemma}[Corollary]
\label{lemma-hilbert-tangent-first-order}
If $Y$ is a closed subscheme of the
projective space $X=\mathbb{P}^n_k$, then the
Zariski tangent space of the Hilbert scheme
$H$ at the point $y$ corresponding to $Y$ is
isomorphic to $H^0(Y,\mathcal{N}_{Y/X})$.
\end{lemma}

Next we consider Situation B. Let $X$ be a
scheme over $k$ and let $L$ be an invertible
sheaf on $X$. We will study the set of
isomorphism classes of invertible sheaves
$L'$ on
$$
X'=X\times D
$$
such that $L'\otimes\mathcal{O}_X\cong L$.
In this case flatness is automatic, because
$L'$ is locally free and $X'$ is flat over
$D$.

\begin{proposition}
\label{proposition-line-bundle-first-order}
Let $X$ be a scheme over $k$, and $L$ an
invertible sheaf on $X$. The set of
isomorphism classes of invertible sheaves
$L'$ on $X\times D$ such that
$L'\otimes\mathcal{O}_X\cong L$ is in natural
one-to-one correspondence with elements of
the group $H^1(X,\mathcal{O}_X)$.
\end{proposition}

\begin{definition}
\label{definition-coherent-first-order}
Let $X$ be a scheme over $k$, and let $F$ be
a coherent sheaf on $X$. We define a
\emph{deformation of $F$ over $D$} to be a
coherent sheaf $F'$ on $X'=X\times D$, flat
over $D$, together with a homomorphism
$F'\to F$ such that the induced map
$$
F'\otimes_D k\longrightarrow F
$$
is an isomorphism. We say that two such
deformations $F'_1\to F$ and $F'_2\to F$ are
\emph{equivalent} if there is an isomorphism
$$
F'_1\xrightarrow{\sim}F'_2
$$
compatible with the given maps to $F$.
\end{definition}

\begin{theorem}
\label{theorem-coherent-first-order}
Let $X$ be a scheme over $k$, and let $F$ be
a coherent sheaf on $X$. The (equivalence
classes of) deformations of $F$ over $D$ are
in natural one-to-one correspondence with
the elements of the group
$\operatorname{Ext}^1_X(F,F)$, where the
zero-element corresponds to the trivial
deformation.
\end{theorem}

\begin{remark}
\label{remark-unmarked-coherent-first-order}
Given $F$ on $X$, we can also pose a
different problem, like the one in (2.6),
namely to classify isomorphism classes of
coherent sheaves $F'$ on $X'$, flat over
$D$, such that $F'\otimes_D k$ is isomorphic
to $F$ (without specifying the isomorphism).
This set need not be the same as the set of
deformations of $F$, but we can explain
their relationship as follows. The group
$\operatorname{Aut}F$ of automorphisms of
$F$ acts on the set of deformations of $F$
by letting $\alpha\in\operatorname{Aut}F$
applied to $f:F'\to F$ be
$\alpha f:F'\to F$. Now let $f:F'\to F$ and
$g:F''\to F$ be two deformations of $F$. One
sees easily that $F'$ and $F''$ are
isomorphic as sheaves on $X'$ if and only if
there exists an
$\alpha\in\operatorname{Aut}F$ such that
$\alpha f$ and $g$ are equivalent as
deformations of $F$. Thus the set of $F'$'s
up to isomorphism as sheaves on $X'$ is the
orbit space of $\operatorname{Ext}^1_X(F,F)$
under the action of $\operatorname{Aut}F$.
This kind of subtle distinction will play an
important role in questions of
pro-representability (Chapter 3).
\end{remark}

\begin{lemma}[Corollary]
\label{lemma-vector-bundle-first-order}
If $E$ is a vector bundle over $X$, then the
deformations of $E$ over $D$ are in natural
one-to-one correspondence with the elements
of $H^1(X,\mathcal{E}nd E)$, where
$\mathcal{E}nd E=\mathcal{H}om(E,E)$ is the
sheaf of endomorphisms of $E$. The trivial
deformation corresponds to the zero element.
\end{lemma}



\section{Deformations of schemes}
\label{section-deformations-of-schemes}

\noindent
The following definition is from
[lucas-defos], which in turn comes from
\cite{Sernesi-deformations}

\begin{definition}
\label{definition-deformation}
Let $X$ be an algebraic $\mathbb{C}$-scheme.
\begin{enumerate}
\item A {\it deformation} of $X$ is a
Cartesian diagram $\xi$
$$
\xymatrix{
X\ar[r]\ar[d]&\mathcal{X}\ar[d]^{\pi}\\
\Spec \mathbb{C}\ar[r]^s&S
}
$$
where $\pi$ is a flat surjective morphism of
algebraic $\mathbb{C}$-schemes and
$S$ is connected. (Recall that flatness
accounts for ``continuity''.)
\item A {\it local deformation} of $X$ is a
deformation $\xi$ where
$S=\Spec A$ for $A$ a noetherian local
$\mathbb{C}$ algebra with residue
field $\mathbb{C}$.
\item An {\it infinitesimal deformation of
$X$} is a local deformation with $A$
an artinian local $\mathbb{C}$-algebra with
residue field $\mathbb{C}$. $X$ is
called {\it rigid} if all infinitesimal
deformations are trivial.
\item An {\it inifinitesimal deformation of
order $n$} is an infinitesimal
deformation when $S=\Spec
(\mathbb{C}[t]/(t^{n+1})$.
\end{enumerate}
\end{definition}


\section{The deformation functor
\texorpdfstring
{$\operatorname{Def}_X$}{Def X}}
\label{section-deformation-functor}

\noindent
The correct object parametrizing abstract
infinitesimal deformations is initially a
functor, not a scheme.

Fix a field $k$ and a $k$-scheme $X$. Let
$\mathbf{Art}_k$ denote the category of
local Artinian $k$-algebras whose residue
field is $k$. Its morphisms are local
$k$-algebra homomorphisms.

\begin{definition}
\label{definition-deformation-artin}
Let $A\in\mathbf{Art}_k$. A
\emph{deformation of $X$ over $A$} is a
flat morphism
$$
\pi:\mathcal{X}_A\longrightarrow
\operatorname{Spec}A
$$
together with an isomorphism, called the
marking of the special fiber,
$$
\iota:\mathcal{X}_A
\times_{\operatorname{Spec}A}
\operatorname{Spec}k\xrightarrow{\sim}X.
$$
An isomorphism between two such deformations
is an isomorphism of $A$-schemes compatible
with their markings.
\end{definition}

\begin{definition}
\label{definition-def-x}
The \emph{deformation functor of $X$} is
defined as follows. For each
$A\in\mathbf{Art}_k$, let $\mathcal{D}_X(A)$
be the set of pairs $(\pi,\iota)$ in the
previous definition. Its value at $A$ is
$$
\operatorname{Def}_X(A)
=
\mathcal{D}_X(A)/{\cong}.
$$
These sets form the functor
$$
\begin{aligned}
\operatorname{Def}_X:
\mathbf{Art}_k&\longrightarrow
\mathbf{Sets},\\
A&\longmapsto
\operatorname{Def}_X(A).
\end{aligned}
$$
If $A\to B$ is a homomorphism, the
corresponding map
$\operatorname{Def}_X(A)\to
\operatorname{Def}_X(B)$ is given by base
change:
$$
\mathcal{X}_A\longmapsto
\mathcal{X}_A\times_{\operatorname{Spec}A}
\operatorname{Spec}B.
$$
\end{definition}

\noindent
A richer version assigns to $A$ the
\emph{groupoid} of deformations and their
isomorphisms. Passing to the set of
isomorphism classes forgets the automorphisms
of each deformation. Infinitesimal
automorphisms can prevent the resulting
set-valued functor from being representable.

\medskip\noindent
\textbf{Comparison with the Hilbert
functor.}
If $X\subset Y$ is a fixed embedding, the
local Hilbert functor parametrizes flat
deformations
$$
\mathcal{X}_A\subset Y\times
\operatorname{Spec}A
$$
of $X$ \emph{inside the fixed ambient
scheme $Y$}. There is a natural
transformation
$$
\operatorname{Def}_{X\subset Y}
\longrightarrow
\operatorname{Def}_X
$$
which forgets the embedding. The functor on
the left describes embedded deformations;
the one on the right describes abstract
deformations.

In our Stanley--Reisner problem, perturbing
the homogeneous generators of
$I_X\subset k[x_0,\ldots,x_7]$ initially
studies embedded deformations of
$X\subset \mathbb{P}^7$. It belongs first
to the local Hilbert problem, although
one can subsequently forget the embedding
and obtain an abstract deformation.

\medskip
\noindent
\textbf{Universal and miniversal formal
deformations.}

\noindent
Let $R$ be a complete local Noetherian
$k$-algebra with residue field $k$. It
defines a functor. Let $\mathbf{C}_k$ denote
the category of such algebras and local
$k$-algebra homomorphisms. We regard every
object of $\mathbf{Art}_k$ as an object of
$\mathbf{C}_k$. Then
$$
h_R(A)=\operatorname{Hom}_{\mathbf{C}_k}
(R,A).
$$
Geometrically, this is the functor associated
to the formal scheme $\operatorname{Spf}R$.

If there is a natural isomorphism
$$
h_R\xrightarrow{\sim}
\operatorname{Def}_X,
$$
then $\operatorname{Def}_X$ is
\emph{prorepresentable}. Equivalently,
$\operatorname{Spf}R$ carries a universal
formal deformation: every infinitesimal
deformation is obtained by pullback along a
unique homomorphism $R\to A$.

In general this is too much to expect. One
therefore asks for a natural transformation
$$
h_R\longrightarrow\operatorname{Def}_X
$$
which is formally smooth and induces an
isomorphism on tangent spaces. Formal
smoothness means that compatible data over
the two rings in a small extension can be
lifted to $h_R$ over the larger ring.
Concretely, for every small extension
$B\longrightarrow A$, the map
$$
h_R(B)\longrightarrow
h_R(A)\times_{\operatorname{Def}_X(A)}
\operatorname{Def}_X(B)
$$
is required to be surjective.

The tangent space of the deformation functor
is
$$
T\operatorname{Def}_X=
\operatorname{Def}_X
\big(k[\varepsilon]/(\varepsilon^2)\big).
$$
Thus the minimality condition says that
$$
T_o\operatorname{Spf}R
\xrightarrow{\sim}
T\operatorname{Def}_X.
$$

\begin{definition}
\label{definition-miniversal-deformation}
A formal deformation over
$\operatorname{Spf}R$ satisfying the two
conditions above is called a
\emph{miniversal deformation}. The formal
scheme $\operatorname{Spf}R$ is called a
\emph{miniversal deformation space}, a
\emph{semiuniversal deformation space}, or
a \emph{hull} of $\operatorname{Def}_X$.
\end{definition}

\noindent
The crucial distinction is that
$\operatorname{Def}_X$ is the functor,
whereas $\operatorname{Spf}R$, together
with the natural transformation above, is
a miniversal space for it.
A miniversal family contains all
infinitesimal deformation directions and has
no redundant tangent parameters, but the map
producing a given deformation need not be
unique. If it were unique for every Artinian
base, the family would be universal and the
functor would be prorepresentable.

Schlessinger's theorem gives conditions under
which a functor on $\mathbf{Art}_k$ has such
a hull, and stronger conditions under which
it is prorepresentable. Once its lifting
conditions hold, finite-dimensionality of
the tangent space is the key finiteness
condition for the existence of a hull.
Proper schemes in standard finite-type
settings have this finiteness. Infinitesimal
automorphisms can prevent a hull from being
universal.

The cotangent complex gives the standard
tangent and obstruction spaces
$$
T^1_X=\operatorname{Ext}^1
(\mathbb{L}_{X/k},\mathcal{O}_X),
\qquad
T^2_X=\operatorname{Ext}^2
(\mathbb{L}_{X/k},\mathcal{O}_X).
$$
The second group is an obstruction space:
obstructions to lifting lie in it, although
not every element must occur as an
obstruction.

Suppose that $T^1_X$ is finite-dimensional
and that a hull exists. A minimal
presentation of its complete local ring has
the form
$$
R\cong k[[t_1,\ldots,t_r]]/J,
\qquad
J\subset (t_1,\ldots,t_r)^2,
\qquad
r=\dim_kT^1_X.
$$
The ideal $J$ records the higher-order
lifting conditions detected by the
obstruction theory. Its lack of linear terms
expresses the tangent-space minimality of the
presentation.

Finally, a formal one-parameter deformation
is only a formal arc in this miniversal
space:
$$
\operatorname{Spf}k[[s]]
\longrightarrow
\operatorname{Spf}R.
$$
Concretely, an arc through the origin is a morphism
$$
\gamma:\operatorname{Spf}k[[t]]
\longrightarrow
\operatorname{Spf}R
$$
sending the closed point $t=0$ to the point representing
$X$. Since morphisms of affine formal schemes reverse the
arrows on rings, this is equivalently a continuous local
$k$-algebra homomorphism
$$
\gamma^\#:R\longrightarrow k[[t]].
$$
If
$$
R\cong k[[x_1,\ldots,x_r]]/J,
$$
then specifying the arc amounts to choosing power series
$$
x_i\longmapsto x_i(t)\in t\,k[[t]]
$$
which satisfy the equations in $J$. Thus the arc selects a
one-dimensional formal path through the possibly much
larger deformation space. Pulling the miniversal family
back along $\gamma$ produces a formal one-parameter
deformation over $\operatorname{Spf}k[[t]]$. In particular,
the arc is a device for passing from the full deformation
space to a one-parameter deformation; it is not a section
of the total family.

In our problem the Maurer--Cartan
construction seeks a compatible deformation
along such an arc. If that formal deformation
is algebraized, one can then ask whether its
generic fiber is smooth. The construction
does not by itself produce the whole
miniversal deformation space.

\section{The Maurer--Cartan equation in complex geometry}
\label{section-maurer-cartan-complex-geometry}

\noindent
Let $X$ be a complex manifold and let $M$ be its
underlying real manifold. The complex structure is
equivalently encoded by the splitting
$$
T_{\mathbb C}M=T^{1,0}X\oplus T^{0,1}X.
$$
Thus, to deform the complex structure we may deform
the subbundle $T^{0,1}X$ inside $T_{\mathbb C}M$.

A sufficiently small deformation which stays transverse
to $T^{1,0}X$ is the graph of a bundle map
$$
\phi:T^{0,1}X\longrightarrow T^{1,0}X.
$$
Equivalently,
$$
\phi\in A^{0,1}(T^{1,0}X),
$$
since
$$
\operatorname{Hom}(T^{0,1}X,T^{1,0}X)
\cong (T^{0,1}X)^*\otimes T^{1,0}X.
$$
The deformed antiholomorphic tangent bundle is
$$
T^{0,1}_\phi
=\{v+\phi(v):v\in T^{0,1}X\}.
$$

Why does such a map really describe a deformation of the
subbundle? This is simply fiberwise linear algebra. Suppose
$$
E=V\oplus W
$$
and let $W'\subset E$ be another subspace of the same
dimension as $W$, transverse to $V$. Then the projection
$$
\pi_W:E\longrightarrow W
$$
restricts to an isomorphism
$$
\pi_W|_{W'}:W'\xrightarrow{\sim}W.
$$
Therefore for every $w\in W$ there is a unique vector of
$W'$ projecting to $w$. It has the form
$$
w+\phi(w)
$$
for a unique $\phi(w)\in V$. Thus there is a unique linear
map
$$
\phi:W\longrightarrow V
$$
such that
$$
W'=\{w+\phi(w):w\in W\}=\operatorname{graph}(\phi).
$$

Applying this at each $x\in X$ with
$$
E=T_{\mathbb C,x}M,\qquad
V=T^{1,0}_xX,\qquad
W=T^{0,1}_xX,
$$
shows that every nearby $n$-plane transverse to
$T^{1,0}_xX$ is uniquely the graph of a map
$$
\phi_x:T^{0,1}_xX\longrightarrow T^{1,0}_xX.
$$
If these planes vary smoothly with $x$, the maps $\phi_x$
assemble into the bundle map $\phi$ above. In other words,
$$
\operatorname{Hom}(T^{0,1}_xX,T^{1,0}_xX)
$$
gives local coordinates near $T^{0,1}_xX$ in the
Grassmannian of $n$-planes in $T_{\mathbb C,x}M$. Thus
$\phi$ is literally the slope of the deformed subbundle
relative to the original splitting.

Here ``sufficiently small'' just guarantees that
$T^{0,1}_\phi$ and its complex conjugate remain
transverse, so that they define an almost complex
structure. For example, after choosing a Hermitian
metric and assuming $X$ compact, it is enough to take
$\phi$ sufficiently small in the $C^0$ norm. Stronger
topologies such as $C^k$, H\"older or Sobolev topologies
are useful when studying the differential equation below.

The remaining question is whether this almost complex
structure is integrable. By the Newlander--Nirenberg
theorem, integrability is equivalent to
$$
[T^{0,1}_\phi,T^{0,1}_\phi]
\subseteq T^{0,1}_\phi.
$$

To compute this condition, choose local holomorphic
coordinates $z^1,\ldots,z^n$ for the original complex
structure and write
$$
\phi=
\phi^k_{\bar i}\,d\bar z^i\otimes
\frac{\partial}{\partial z^k}.
$$
Then $T^{0,1}_\phi$ is locally spanned by
$$
e_{\bar i}
=\frac{\partial}{\partial\bar z^i}
+\phi^k_{\bar i}\frac{\partial}{\partial z^k}.
$$
A direct computation gives
$$
[e_{\bar i},e_{\bar j}]
=
\left(
\frac{\partial\phi^k_{\bar j}}{\partial\bar z^i}
-\frac{\partial\phi^k_{\bar i}}{\partial\bar z^j}
+\phi^\ell_{\bar i}
\frac{\partial\phi^k_{\bar j}}{\partial z^\ell}
-\phi^\ell_{\bar j}
\frac{\partial\phi^k_{\bar i}}{\partial z^\ell}
\right)
\frac{\partial}{\partial z^k}.
$$
The first two terms are the components of
$\bar\partial\phi$, while the last two are the
components of $\frac12[\phi,\phi]$, where the bracket
is the natural extension of the Lie bracket of vector
fields to $T^{1,0}X$-valued forms. Hence the
integrability condition is precisely
$$
\boxed{
\bar\partial\phi+\frac12[\phi,\phi]=0.
}
$$
This is the \emph{Maurer--Cartan equation}.

Thus the differential graded Lie algebra
$$
A^{0,\bullet}(T^{1,0}X),
\qquad d=\bar\partial,
$$
with this bracket naturally appears when one studies
deformations of complex structures: its
Maurer--Cartan elements are precisely the small
integrable deformations written in these graph
coordinates.

Up to this point, $\phi$ has represented one single
nearby complex structure: its graph is the deformed
subbundle $T^{0,1}_\phi$. There has been no deformation
parameter $t$ yet.

To discuss first order, second order and lifting, we now
make an additional choice: consider a one-parameter family
of deformations passing through the original complex
structure,
$$
t\longmapsto \phi_t,
\qquad
\phi_0=0.
$$
Thus we regard the family as a curve,
that is, as vector-valued 
map of a real parameter $t$, 
with values in the
vector space
$$
A^{0,1}(T^{1,0}X)
=\Gamma\!\left(\operatorname{Hom}
(T^{0,1}X,T^{1,0}X)\right).
$$
If this family depends analytically on $t$, we may Taylor
expand it at $t=0$. Absorbing the usual factorials into the
coefficients, we write
$$
\phi(t)=t\phi_1+t^2\phi_2+t^3\phi_3+\cdots.
$$
In formal deformation theory we instead take such a formal
power series as the object from the start. The condition
$\phi(0)=0$ explains why there is no constant term: at
$t=0$ we recover the original complex structure.

For this one-parameter formal deformation, the complete
local base ring is
$$
\mathbb C[[t]]
=\varprojlim_n \mathbb C[t]/(t^{n+1}).
$$
Thus a formal deformation over $\mathbb C[[t]]$ means a
compatible system of finite-order deformations over
$$
\operatorname{Spec}\big(\mathbb C[t]/(t^{n+1})\big),
$$
or equivalently a deformation over the formal base
$\operatorname{Spf}\mathbb C[[t]]$. The series $\phi(t)$
above is therefore naturally a power series in the
coordinate $t$ of this formal base.

Here ``first order'' means the coefficient of $t^1$ after
expanding the Maurer--Cartan equation as a formal power
series in $t$. Indeed,
$$
\bar\partial\phi(t)
=t\bar\partial\phi_1+t^2\bar\partial\phi_2+\cdots,
$$
while bilinearity of the bracket gives
$$
[\phi(t),\phi(t)]
=t^2[\phi_1,\phi_1]
+2t^3[\phi_1,\phi_2]+\cdots.
$$
Thus
$$
t\bar\partial\phi_1
+t^2\left(
\bar\partial\phi_2+\frac12[\phi_1,\phi_1]
\right)
+\cdots=0.
$$
A formal power series is zero if and only if each of its
coefficients is zero. Hence
$$
\boxed{\text{order }t:\quad \bar\partial\phi_1=0}
$$
and
$$
\boxed{
\text{order }t^2:\quad
\bar\partial\phi_2+\frac12[\phi_1,\phi_1]=0.
}
$$
Equivalently, working to first order means working modulo
$t^2$. Then
$$
\phi(t)\equiv t\phi_1\pmod{t^2},
$$
and the bracket term disappears because it is already of
order $t^2$. This is why first-order deformation theory is
linear, whereas the first obstruction appears at second
order.

Thus the bracket produces the first quadratic
obstruction: {\bf the first-order direction $\phi_1$ lifts
to second order only if
$\frac12[\phi_1,\phi_1]$ is
$\bar\partial$-exact}.
The corresponding obstruction is therefore
the cohomology class
$$
\left[\frac{1}{2}[\phi_1,\phi_1]\right]
\in
H^{0,2}(X,T^{1,0}X).
$$
If that class is zero, then we have
Maurer-Cartan and we can choose $\phi_2$
to lift $\phi_1$ to second order.
But if this class is nonzero, we are stuck.


After quotienting first-order deformations by
infinitesimal changes of coordinates, their classes lie
in
$$
H^{0,1}(X,T^{1,0}X)\cong H^1(X,T_X),
$$
which is exactly the space that appears in the
Kodaira--Spencer map below.

The isomorphism above is simply the Dolbeault theorem for
the holomorphic vector bundle $T_X$. The sheaf $T_X$ of
holomorphic vector fields has the Dolbeault resolution
$$
0\longrightarrow T_X
\longrightarrow \mathcal A^{0,0}(T^{1,0}X)
\xrightarrow{\bar\partial}
\mathcal A^{0,1}(T^{1,0}X)
\xrightarrow{\bar\partial}
\mathcal A^{0,2}(T^{1,0}X)
\xrightarrow{\bar\partial}\cdots .
$$
This is an exact resolution of $T_X$, and the sheaves
$\mathcal A^{0,q}(T^{1,0}X)$ are fine, hence acyclic for
global sections. Therefore the sheaf cohomology of $T_X$
is computed by the cohomology of the complex of global
sections
$$
A^{0,0}(T^{1,0}X)
\xrightarrow{\bar\partial}
A^{0,1}(T^{1,0}X)
\xrightarrow{\bar\partial}
A^{0,2}(T^{1,0}X)
\longrightarrow\cdots .
$$
Consequently,
$$
H^q(X,T_X)\cong H^{0,q}_{\bar\partial}
(X,T^{1,0}X),
$$
and in particular
$$
H^1(X,T_X)\cong H^{0,1}_{\bar\partial}
(X,T^{1,0}X).
$$
So the same quotient that appears naturally from the
Maurer--Cartan linearization---$\bar\partial$-closed
first-order directions modulo $\bar\partial$-exact
infinitesimal coordinate changes---is precisely the usual
sheaf cohomology group $H^1(X,T_X)$.

\medskip\noindent
The point to retain is that the Maurer--Cartan equation
is not an extra formal gadget imposed on deformation
theory. Already for ordinary complex structures it is
simply the integrability condition for the deformed
object. Its linear term describes first-order
deformations, its bracket gives the quadratic
obstruction, and the full equation organizes lifting to
all orders.

\section{Kodaira-Spencer map}
\label{section-kodaira-spencer-map}

\noindent
Let $\mathcal{X} \to B$ be a
deformation of $X$.
Kodaira-Spencer map is
$$
T_pB \to H^1(X,T_X)
$$
constructed as follows.

\begin{enumerate}
\item Since $\pi$ is a submersion,
we can take submersion charts all along
the fiber $X = \pi^{-1}(p)$. This gives
a cover $\mathcal{U}_i\simeq U_i \times V_i$ 
where
$V_i$ is the part of the chart taken
using the coordinates of $B$
(this is just inverse function
theorem in its level-set chart theorem
incarnation).

\item In the intersections we have
transition functions:
$$
\xymatrix{
\mathcal{U}_i \cap \mathcal{U}_j
\ar[r]^{\varphi_i}& 
U_i \times V \ar[d]^{\varphi_{ij}}\\
\mathcal{U}_i \cap\mathcal{U}_j
\ar[r]_{\varphi_j}&
U_j \times V
}
$$
where $V:= V_i \cap V_j$ may thought
of as a neighbourhood of $B$
which must contain $p$.
This means that we can apply the
Regular Level-set theorem using the
same coordinates for the $V$-part
of the trivializations.

Then the transition functions
$\varphi_{ij}$ look like this:
$(\psi_{ij}(x,t),t)$ 
where $x$ is a variable
in $X = \pi^{-1}(p)$, 
$t$ is a variable in $B$
and $\psi_{ij}$ is a transition map from the
original atlas of $\mathcal{X}$.

\item Taking a tangent vector $v$ 
at $p \in B$ we get an 
integral curve $\gamma$.
Differentiating we obtain a vector field
on $U_i \cap U_j$:
$$
\theta_{ij}=\frac{d}{dt}\Big|_{t=0}
\psi_{ij}(x,\gamma(t)).
$$
\item Since the $\psi_{ij}$
satisfy the cocycle condition
(by the definition of the atlas
transition maps 
$\psi_{ij}=\psi_i^{-1}\circ \psi_j^{-1}$),
so do the $\theta_{ij}$, constituting
a cocycle in $H^1(T_X)$.
\end{enumerate}


\section{Kodaira-Spencer correspondence}
\label{section-kodaia-spencer-correspondence}


\medskip\noindent
After a few months of the previous text in
this section
I think I have understood the construction of
Kodaira-Spencer map for embedded
deformations,
which is \cite[Proposition
3.2.1(ii)]{Sernesi-deformations}.

But I will actually briefly outline a
particular case
of this, which is easier: \cite[Examples
3.2.4(ii)]{Sernesi-deformations}.
This the case when $X \subset Y$ is a Cartier
divisor.
Then $X$ is cut out locally by a set of
functions $f_i$
defined on an affine cover $U_i$ of $Y$.
Then we have a line bundle defined by the
transition functions $f_{ij}=f_i/f_j$.
This is the line bundle $\mathcal{O}_Y(X)$,
which
by the adjunction formula is also the normal
bundle $N_{X/Y}$
of $X$ in $Y$.

Now consider a first order deformation
$$
\xymatrix{
X\ar[r]\ar[d]
&\mathcal{X}\subset Y\times
B_\varepsilon\ar[d]^\pi\\
0\ar[r]
&B_\varepsilon
}
$$

The great thing about first order
deformations
is that the structure sheaf of $\mathcal{X}$
on an affine open $U_i$ of $Y$
looks like the tensor product
$\mathcal{O}_{U_i}\otimes
\mathbb{C}[\varepsilon]/(\varepsilon^2)$
(\href{https://stacks.math.columbia.edu/tag/%
01I4}{01I4}).
By definition of tensor product this just
says
that sections look like $f+\varepsilon g$
for $f,g \in \Gamma(\mathcal{O}_{U_i})$.

If $\mathcal{X}$ is a Cartier divisor on $Y
\times B_\varepsilon$
(it is; but I'm missing a formal proof;
the intuition is that an ``infinitesimal
thickening'' $\mathcal{X}$
of $X$ has the same underlying topological
space
but different scheme structure;
\href{https://stacks.math.columbia.edu/tag/%
01JL}{01JL}
\href{https://stacks.math.columbia.edu/tag/%
01I4}{01I4}),
it is given as the vanishing locus
of some local sections $f_i+\varepsilon g_i$.
(Perhaps the $f_i$ coincide with the
distinguished
section of $X$.)
The line bundle associated to $\mathcal{X}$
is given by the cocycle
$$
F_{ij}=f_{ij}+\varepsilon g_{ij}
$$
for $g_{ij}=\frac{g_i-f_{ij}g_j}{f_j}$.
Indeed, note that
$\frac{1}{f_j+\varepsilon
g_j}=\frac{1}{f_j}-\varepsilon
\frac{g_j}{f^2_j}$
and compute $F_{ij}=(f_i+\varepsilon
g_j)/(f_j+\varepsilon g_j)$.

Then
$$
f_i+\varepsilon g_i=(f_{ij}+\varepsilon
g_{ij})(f_j+\varepsilon g_j)
$$
gives
$$
g_i=f_{ij}g_j+g_{ij}f_j
$$
and since $f_j$ vanishes along $X$,
we are left with $g_i=f_{ij}g_j$
which is how a section of the sheaf $N_{X/Y}$
is constructed by definition of sheaf.


\begin{example}
\label{example-modul-space-nonsi-riema}
\begin{reference}
\cite[Example 3.1]{continued-fractions}
\end{reference}
Fix $g \geq 2$. The dimension of the
deformation space of a nonsingular
projective curve $X$ is $3g-3$. This is ``the
dimension of the moduli space of
curves of genus $g$''.

We can compute this number by Riemann-Roch
formula on the bundle $-K_X$.
Indeed, since $X$ is a curve,
$\Omega_X^1=K_X$ and thus $-K_X=T_X$. We get
$$
\begin{aligned}
h^0(-K)-h^0(K-(-K))&=\text{deg}(-K)-g+1\\
&=-2g+2-g+1\qquad \substack{\text{degree
additive and}\\ \text{deg}K=2g-2}\\
&=-3g+3
\end{aligned}
$$
Now by Serre duality
$h^0(K-(-K))=h^1(-K)=h^1(T_X)$, and it turns
out that a
Riemann surface of genus $g \geq 2$ has no
holomorphic vector fields, so that
$h^0(-K)=h^0(T_X)=0$.
\end{example}

\begin{exercise}
\label{exercise-deformations-of-curves-in-K3}
Let $C$ be a smooth genus $g$ curve which can
be embedded in a K3 surface  $M$,
and $X$ the family of all deformations of $C$
in $M$.
\begin{enumerate}
\item Prove that $\dim X \leq  g$.
\item Let $\mathcal{X}_g$ be the space of all
curves of genus $g$ (smooth?)
which can be possibly embedded to a K3
surface. Prove that each irreducible
component $Z$ of $\mathcal{X}_g$ satisfies
$\dim_\mathbb{C} \leq  g+19$. Deduce
that there exists a compact complex curve
which cannot be embedded in a K3
surface.
\end{enumerate}
\end{exercise}

\begin{exercise}
\label{exercise-}
Let $C$ be a smooth curve embedded in a K3
surface $X$. Show that the dimension
of the deformation space of $C$ is $\leq g$.
\end{exercise}

\begin{proof}
% \begin{enumerate}
% \item
The deformation space of a variety is the
space of isomorphism classes of
deformations as explained above. It turns out
that there is a way to associate
1-cocyles of the tangent sheaf to
deformations, so that in fact the deformation
space $\text{Def}_1$ is isomorphic to
$H^{1}(X,\mathcal{T}_X)$ for any variety
$X$.

For our curve $C$ we thus know that the
dimension of the space of deformations
(deformations not necessarily contained in
$X$) is $h^1(\mathcal{T}_C)$.  The
family of deformations of $C$ that are
contained in $X$ is the Hilbert scheme of
curves with fixed Hilbert polynomial $P(t)$
after quotienting by $\mathbb{C}s$,
where $s$ is the section whose vanishing
locus is  $C$.
This says that the number we are looking for
is $h^{0}(X,\mathcal{O}(C))-1$.

Before showing that in fact
$h^0(X,\mathcal{O}(C))=1+g$,
let's explain why this Hilbert scheme
construction actually works.
Intuitively, any curve lying in the same
component
as $C$ would be deformation-equivalent to $C$
--- at least we know that it's Hilbert
polynomial is the same.
But it's not clear what's the deformation
involving these two
(this question should be addressed!).
Further, this doesn't explain Vladimiro's
suggesion of ``taking quotient by
$\mathbb{C}s$''.
To me it looks like I would just compute the
dimension
of the tangent space of $\text{Hilb}^P(M)$ at
$C$.

Now I will show that
$h^0(X,\mathcal{O}(C))=1+g$ (see
\cite[Lemma 1.2.1, Remark 1.2.2]{huk}).

Consider the ideal sheaf exact sequence
twisted by
$\mathcal{O}(C)$:
$$
\xymatrix{
0\ar[r]&\mathcal{O}_X\ar[r]&\mathcal{O}(C)
\ar[r]
&\mathcal{O}_X(C) \otimes
\mathcal{O}_C=\mathcal{O}(C)|_{C}\ar[r]&0
}
$$
By $X$ being a K3 we know that
$H^{1}(\mathcal{O}_X)=0$, so that we have the
short exact sequence in cohomology
$$
\xymatrix{
0\ar[r]&H^{0}(\mathcal{O}_X)\ar[r]&H^{0}
(\mathcal{O}(C))\ar[r]
&H^{0}(\mathcal{O}(C)|_{C})\ar[r]&0
}
$$
so that
$h^0(\mathcal{O}(C))=1+h^1(\mathcal{O}(C)
|_{C})$.
A version of adjunction formula says
$\omega_C \cong (K_X
\otimes\mathcal{O}(C)|_{C}$, and using that
$K_X=\mathcal{O}_X$ we obtain
$h^0(\mathcal{O}(C)|_{C})=h^0(\omega_C)=g$.

To

\medskip\noindent
For the record I put other thoughts I went
through in solving this exercise.

Recall from \cite[p. 146]{gri} that the
normal bundle
$\mathcal{N}$ of a hypersurface of a smooth
variety $X$ satisfies
$\mathcal{N}^\vee\cong\mathcal{O}_X(-C)|_{C}
$.
Taking duals we get that
$\mathcal{N}\cong\mathcal{O}_X(C)|_{C}$.

By adjunction formula
$2g-2=(\mathcal{O}(C),\mathcal{O}(C))$.
Applying Riemann-Roch
to $\mathcal{O}(C)$ (which is by definition
the dual of the ideal sheaf of $C$,
which is a line bundle on $X$), we obtain
that
$\chi(\mathcal{O}(C))=2+\frac{1}{2}
(\mathcal{O}(C),\mathcal{O}(C))$.
Then
$\chi(\mathcal{O}(C))=g+1$.

Now recall that
$\chi(\mathcal{O}(C))=h^0(\mathcal{O}(C)
-h^1(\mathcal{O}(C))+h^2(\mathcal{O}(C))$. By
Serre duality and $X$ being a K3
surface we see that $h^2(\mathcal{O}(C))\cong
h^{0}(\mathcal{O}(-C))$, which is
the ideal sheaf of $C$. Any section of such a
sheaf would vanish along $C$, and
since $X$ is compact we conclude there cannot
be any such section.

Now we show that also $h^1(\mathcal{O}(C))=0$
to conclude that
$h^0(\mathcal{O}(C))=g+1$.

We thus conclude that $h^0(\mathcal{O}(C))$




Since $C$ is smooth we can use the normal
exact sequence
$$
\xymatrix{
0\ar[r]&\mathcal{T}_C\ar[r]&\mathcal{T}
_X|_{C}\ar[r]&\mathcal{N}\ar[r]&0
}
$$
Taking Euler characteristic we see that
$\chi(\mathcal{T}_C)+\chi(\mathcal{N})=
\chi(\mathcal{T}_X|_{C})$.


This means that we should be done once we
compute $\chi(\mathcal{T}_X|_{C})$.
For this we can use Riemman-Roch formula for
coherent sheaves on a curve, which
tells us that
$$
\text{deg}(\mathcal{T}_X|_{C})=
\chi(\mathcal{T}_X|_{C})
-\text{rk}(\mathcal{T}_X|_{C})
\cdot\chi(\mathcal{O}_C)
$$
But of course we know that since $C$ is a
curve, by Serre duality we get that
$\chi(\mathcal{O}_C)=1-g$. (Indeed:
$h^1(\mathcal{O}_C)=h^0(\Omega^1_C):=p_a(C)
$.)
And now the question is what is
the degree. Apparently this is just the first
Chern class $c_1$ of the
restricted tangent bundle. So what is it? And
what is $h^0(\mathcal{T}_C)$ if
the genus is 0 or 1, and that's it.

I
know that $h^0(\mathcal{T}_X)=0$ and
$h^1(\mathcal{T}_X)=20$ by $X$ being a K3
surface, but I'm not sure what happens when
we restrict to $C$.

If $g \geq 2$ we know that
$h^0(\mathcal{T}_C)=0$, so that
$\boxed{\chi(\mathcal{T}_C)=h^1(\mathcal{T}
_C)}$.
\bigskip
To compute the
latter Euler characteristic, which is given
by definition by
$\chi(\mathcal{T}_X|_{C})=h^0(\mathcal{T}
_X|_{C})-h^1(\mathcal{T}_X|_{C})$,
we
first note that $h^0(\mathcal{T}_X|_{C})=0$,
because this is the dimension of
the global holomorphic vector fields on $X$
restricted to $C$, which is constant along
$C$ since $X$ is
smooth. And then this number is actually zero
by Hodge numbers of a K3 surface.

and taking cohomology long exact sequence we
obtain
$$
\xymatrix{
\cdots\ar[r]&H^{0}(T_X)\ar[r]&H^{0}
(\mathcal{N})\ar[r]&H^{1}(T_C)\ar[r]&\\
H^{1}(T_X)\ar[r]&H^{1}(\mathcal{N})\ar[r]&
\cdots
}
$$
\medskip\noindent
Or is it $h^1(\mathcal{N})$? See
\cite{HarrMorr}. If this was the case, then
we
can use adjunction formula as above to get
that
$2g-2=(\mathcal{N},\mathcal{N})=\text{deg}_C
\mathcal{N}$. Then we may find
$h^0(\mathcal{N})$ via Riemann-Roch:
$$
h^0(\mathcal{N})-h^0(K_C-\mathcal{N})=
\text{deg}\mathcal{N}-g+1
$$
Note that
$h^0(N_C-\mathcal{N})=h^1(-K_C+N+K_C)=
h^1(\mathcal{N})$
via Serre
duality, and by Riemann-Roch on a surface as
above we see that
$$
g+1=\chi(\mathcal{N})=h^0(\mathcal{N})-
h^1(\mathcal{N})
\implies
h^1(\mathcal{N})=-g-1+h^0(\mathcal{N})
$$
so that
$$
\begin{aligned}
h^0(\mathcal{N})-(-g-1+h^0(\mathcal{N}))&=
\text{deg}\mathcal{N}-g+1\\
\implies \text{oops! I lost
$h^0(\mathcal{N})$ in this operation…}
\end{aligned}
$$
Maybe if I just use normal exact sequence and
realise that
$$
h^0(\mathcal{N})=h^0(\mathcal{T}_X|_{C})-
h^0(\mathcal{T}_C)
$$
I know that $h^0(\mathcal{T}_C)=0$ for $g>1$,
so the question is how to compute
the restricted holomorphic vector fields.
% \end{enumerate}
\end{proof}

\section{Deformations of Stanley-Reisner schemes}
\label{section-sr}

\noindent
My aim in this section is to cover the
minimal
theory needed to apply results in \cite{jan2}
to my deformation problem. The main theorem
I expect to use is the following:

\begin{theorem}
\label{theorem-deformations-threefolds}
\begin{reference}
\cite[Theorem 5.7]{jan2}
\end{reference}
If $\mathcal{K}$ is a 3-dimensional manifold
then
$$
\dim T^1_{\mathbb{P}(K)}
=11d_3+5e_3+3e_4+e_{\geq 5}+c_{\geq
6}+5f_1^{(3)}+2f_1^{(4)}+h^2(\mathcal{K}).
$$
\end{theorem}

\noindent
I need to do some interpretation first.
From \cite[p. 1]{jan3}:

\begin{quotation}
  ``For each $\mathbb{C}$-algebra $A$,
there is a cohomology theory providing
modules $T^i_A$.
However, only $T^1_A$ and $T^2_A$ are
relevant for the deformation
theory of $\operatorname{Spec} A$. The module $T^1_A$
collects the infinitesimal
deformations of $A$ and $T^2_A$ contains the
obstructions
for lifting these deformations to decent
parameter spaces.
Eventually, $\text{Der}_{\mathbb{C}}(A,A)$
will be called $T_A^0$."
\end{quotation}

\cite[p. 3]{jan3}:
\begin{quotation}
``In general, for a finitely generated
$\mathbb{C}$-algebra $A$,
the modules $T_A^i$ allow the following ad
hoc definitions:
Let $P=\mathbb{C}[x_0,…,x_n]$ mapping onto
$A$ so that $A \simeq P/I$
for an ideal $I$. The $T_A^1$ is the cokernel
of the natural map
$\text{Der}_\mathbb{C}(P,P) \to \Hom_P(I,A)$.
Moreovr, […] and we obtain $T^2_A$ as the
cokernel
of the induced map
$\Hom_P(P^m,A)\to\Hom_A(R/R_0,A)$.''
\end{quotation}

\begin{remark}[dani]
\label{remark-dani1}
Let's think more about that definition of
$T^1_A$
as the cokernel of the ``natural map''
$\text{Der}_\mathbb{C}(P,P) \to \Hom_P(I,A)$.

Pick a derivation $\chi \in
\text{Der}_\mathbb{C}(P,P)$.
We want a map $I \to A=P/I$.
Pick an element $a \in I$ and map it to
$\chi(a)$.
And that's it. It might be zero or not.
If the $\chi$ we picked maps any $a \in I$ to
some element in $I$,
then this $\chi$ is in the kernel of the
natural map.
So the image consists of the deformations
coming from infinitesimal coordinate
changes. The module $T^1_A$ is the target
modulo this image, namely the cokernel.
\end{remark}

\noindent
Unfortunately this doesn't say much about the
following crucial fact:
the $T^i_A$ modules are
$\mathbb{Z}^{n+1}$-graded.
But here goes an explanation.
We use the notation that $e_p$ is a ``face
not in the complex'',
which of course corresponds to a generating
monomial $x^p$
of the SR ideal $I$.
It has a degree, of course (given by the
``support'' of $p$).
Now here's a rather sloppy step but:
since the cotangent modules are built upon
$I$,
who is given by $e_p$,
we just notice that the degree of $e_p$
along with a (very natural, too) notion of
degree among relations,
induce a grading in all the cotangent modules
$T_A^i$.
Thus, the grading is morally just the degree
of the monomials,
which is in fact the support of $p$, a vector
$\mathbf{c} \in \mathbb{Z}^{n+1}$.

\medskip\noindent
{\bf The $a-b$ grading.}
Next in line is the decomposition
$\mathbf{c}=\mathbf{a}-\mathbf{b}$
into positive minus negative part. But how
exactly…? 

Let $I_{\mathcal M}$ be a Stanley--Reisner
ideal. Its generators are monomials
corresponding to missing faces of
$\mathcal M$.

A first-order deformation replaces a
forbidden monomial by that monomial plus a
legal monomial. Schematically,
$$
x_b \quad \leadsto \quad x_b + t x_a .
$$

Here $x_b$ is the original Stanley--Reisner
generator. Thus $b$ records the missing face.

The monomial $x_a$ is the correction term.
Its support is a face of the complex, or at
least lies in a face/facet region of the
complex. Thus $a$ records what is inserted
by the deformation.

The multidegree of this deformation is
$$
\deg(x_a)-\deg(x_b)=a-b.
$$

For example, the direction $t_{48}$ sends
the missing triangle $348$ to the monomial
$x_2x_5^2$. Thus
$$
x_3x_4x_8
\quad\leadsto\quad
x_3x_4x_8+t x_2x_5^2.
$$
So in this case
$$
b=348,
\qquad
a=\deg(x_2x_5^2).
$$

In words: $b$ is the forbidden face being
deformed, and $a$ is the legal monomial
inserted as the first-order correction.




\section{Minimal resolutions}
\label{section-minimal-resolutions}

\noindent
In this section I explain what is a minimal
resolution and what are syzygies. The
importance of this is that our main tool to
obtain flatness shall be changesiArtin's
criterion for flatness (\cite[Theorem
3.15]{Graffeo-Alessi}), which requires
computation of the first syzygy module.

\medskip\noindent
If an $S$-module $M$ is free, then it won't
have any syzygies. If $M$ is not free, then
its first syzygy module is the module
generated by the relations which define  $M$.
More formally: consider a set $A$ of
generators of $M$ and let $S^A$ be a free
$S$-module generated by $A$. This module maps
onto $M$, and its kernel is the {\it first
syzygy module}:
$$
\xymatrix{
\text{Syz}^1(M)\ar[r]
&S^A\ar[r]
&M\ar[r]
&0.
}
$$
If $\text{Syz}^1(M)$ is free, it won't have any syzygies. But if it is not free, it will have syzygies of itself (because by Hilbert Basis Theorem it is finitely generated), giving
$$
\xymatrix{
\text{Syz}^2(M)\ar[r]
&S^{A_1}\ar[r]
&\text{Syz}(M)
&0
}
$$
where $S^{A_1}$ is the free $S$-module
generated by a set $A_1$ of generators of $\text{Syz}^1(M)$.

The {\it minimal resolution} of $M$ is
$$
\xymatrix{
&S^{A_k}\ar[r]
&\cdots\ar[r]
&S^{A_2}\ar[r]
&S^{A_1}\ar[r]
&M\ar[r]
&0.
}
$$
This means that the syzygy modules
{\bf do not} appear in the minimal resolution!
They are modules lying in-between each of the arrows
joining the free modules.
The minimal resolution is independent
(up to isomorphism) of the choices of generating
sets of each of the syzygy modules.

\medskip\noindent
Consider
$$
\xymatrix{
0\ar[r]
&S(-1)\ar[r]^{\cdot x_0}
&S\ar[r]
&S/(x_0)\ar[r]
&0
}
$$

In $S(-1)$ constans have degree 1, linear
polynomials have degree 2, and so on. This
means that the map $\cdot x_0$ maps degree-1
objects of $S(-1)$ to degree-1 objects in
$S$; i.e., it is degree-preserving. Its image
is $(x_0)$, so its cokernel is $S/\text{img}
= S/(x_0)$.

\medskip\noindent
{\bf Explanation of M2 notation}

\noindent
Usually we see in Macaulay something like
\[
\begin{array}{cccccc}
      & S      & \xleftarrow{} S^{16} & \xleftarrow{} S^{30} & 
	\xleftarrow{} S^{16} & \xleftarrow{} S \\
      &        &        &        &        &    \\
      & 0      & 1      & 2      & 3      & 4      \\
      & 1      & 16     & 30     & 16     & 1
\end{array}
\]

Here the module is again $S/I$, but it
doesn't appear in Macaulay's output. Instead,
$S/I$ is the cokernel of the first map
$S^{16}\to S$, and the map we had before $S
\to M$ is the composition $S^{16}\to S \to
S/I$.



\section{Artin's criterion for flatness}
\label{section-artins-criterion}

\noindent
The roadmap is
(I recall I couldn't get to a point
of absolute understanding when I last
dived into the details of
how to apply this theorem):
\begin{enumerate}
\item Start with a finitely generated
$R$-module $R/I$ with $I = (f_1,…,f_m)$.
Compute it's minimal resolution.

\item Choose an ideal $I_A$ inside
$R \otimes_\mathbb{C} A := R_A$
generated by elements $F_i$
such that $F_i-f_i \in \mathfrak{m}$
where $m$ is the maximal ideal of $A$.
This means that $F_i$ is a {\it lift} of
$f_i$.

\item Choose lifts $S_j$ of the syzygies
$s_i$ of $R/I$ in the same way.
Ask that the lifted syzygies $S_j$
are syzygies of the lifted complex
(this is to say:
$R^\ell_A \to R^m_A\to R_A \to R_A/I_A$
is exact.)

\item In this construction, when we tensor
by  $\mathbb{C}=A/\mathfrak{m}$
we are literally killing the new variables,
thus coming back to the fiber of
the closed point.

\item Then $A \to R_A/I_A$ is flat
with fiber over the closed point
$R \to R/I$.
\end{enumerate}

\begin{theorem}[Artin's criterion for flatness]
\label{theorem-artin-criterion}
Let $R=\mathbb{C}[x_1,…,x_n]$ be the polynomial ring
in $n$ variables and complex coefficients.
Let also $A$ be a local Artinian algebra of finite type over $\mathbb{C}$.
Let also $I \subset R$ and $I_A \subset R_A=R\otimes A$.
Suppose that there is an exact sequence
\begin{equation}
\label{equation-artin-criterion1}
\xymatrix{
R^\ell\ar[r]
&R^m\ar[r]
&R\ar[r]
&R/I\ar[r]
&0
}
\end{equation}
\noindent
and a complex
\begin{equation}
\label{equation-artin-criterion2}
\xymatrix{
R^\ell_A\ar[r]
&R^m_A\ar[r]
&R_A\ar[r]
&R_A/I_A\ar[r]
&0
}
\end{equation}

\noindent
such that the part
$$
\xymatrix{
R^m_A\ar[r]
&R_A\ar[r]
&R_A/I_A\ar[r]
&0
}
$$
is exact. Then, if
\ref{equation-artin-criterion2}$\otimes_A\mathbb{C}
=$\ref{equation-artin-criterion1},
the module $R_A/I_A$ is $A$-flat.
\end{theorem}

\section{Cotangent obstructions and
the Gr\"unbaum quadrics}
\label{section-cotangent-quadrics}

\noindent
The question is concrete: given a flat
first-order deformation of our equations,
can we add second-order terms and keep it
flat? We shall construct the space in which
the failure of this attempt lives.

\medskip\noindent
\textbf{Equations and their syzygies.}
Work over $k=\mathbb{C}$. Put
$$
S=k[x_1,\ldots,x_n],\qquad
I_0=(f_1,\ldots,f_m),\qquad A=S/I_0.
$$
The generators are homogeneous, with
$d_i=\deg f_i$. Write them as a row,
and introduce their generator module:
$$
F_0=(f_1,\ldots,f_m),\qquad
E=\bigoplus_{i=1}^m S(-d_i).
$$
If $\epsilon_i$ denotes the $i$th basis
vector of $E$, then $F_0(\epsilon_i)=f_i$.
The shift gives $\epsilon_i$ degree $d_i$,
so $F_0:E\to S$ preserves degree. Let
$$
\mathrm{Rel}=\ker(F_0:E\to S).
$$
Choose a graded free module $P$ and a matrix
$R_0:P\to E$ whose columns generate
\emph{all} of $\mathrm{Rel}$. Thus
$$
F_0R_0=0.
$$
Each column records an identity among the
generators, with polynomial coefficients.
The columns need not be independent or
minimal. In the Gr\"unbaum computation,
$m=16$ and a minimal $R_0$ has 30 columns.

Here is the precise flatness statement we
need, in the setting of
Theorem~\ref{theorem-artin-criterion}.
Let $B$ be a local Artinian $k$-algebra
with residue field $k$, and put
$S_B=S\otimes_k B$ and $E_B=E\otimes_k B$.
Choose lifts $F$ of the row $F_0$. Then
$S_B/(F)$ is flat over $B$ if and only if
every syzygy of $F_0$ lifts to a syzygy of
$F$. Equivalently, there is a matrix $R$
reducing to $R_0$ with $FR=0$ in $S_B$.

For justification, apply the
\href{https://stacks.math.columbia.edu/%
tag/051K}{Stacks Project, Lemma 10.101.6
(Tag 051K)}, with the maximal ideal of $B$.
It says that flatness is equivalent to
$$
\operatorname{Tor}^B_1(S_B/(F),k)=0.
$$
Indeed, write $J=(F)$. Tensoring
$0\to J\to S_B\to S_B/J\to0$ with $k$
identifies this Tor group with the kernel
of $J\otimes_B k\to S$. A presentation
$E_B\to J$ shows that this kernel vanishes
exactly when the relations among the $f_i$
come from relations among their lifts.
This also explains why generators of the
syzygy module suffice.

The coefficients of a syzygy are therefore
allowed to change. Requiring $FR_0=0$
would impose an extra condition: flatness
only requires some lift $R$ of $R_0$.

\medskip\noindent
\textbf{Expanding to second order.}
Put $B_r=k[t]/(t^{r+1})$, with $t$ of
degree zero in the polynomial grading.
Over $B_2$ write
$$
\begin{aligned}
F&=F_0+tF_1+t^2F_2,\\
R&=R_0+tR_1+t^2R_2.
\end{aligned}
$$
The rows $F_1,F_2:E\to S$ correct the
defining polynomials. The matrices
$R_1,R_2:P\to E$ correct the coefficients
of the syzygies. All maps preserve degree;
each $F_i$ has size $1\times m$, and each
$R_i$ has the same size as $R_0$.

Multiplication, followed by comparison of
the coefficients of $t$ and $t^2$, gives
\begin{equation}
\label{equation-syzygy-lifts}
\begin{aligned}
F_0R_1+F_1R_0&=0,\\
F_0R_2+F_1R_1+F_2R_0&=0.
\end{aligned}
\end{equation}
The constant coefficient already vanishes
because $F_0R_0=0$. These two identities
are exactly $FR=0$ modulo $t^3$.

Suppose the first identity has been solved.
Then $F_1,R_1$ give a flat deformation over
$B_1$. At the next step the known residual
is $F_1R_1$. We must find $F_2,R_2$ such
that
$$
F_1R_1=-F_0R_2-F_2R_0.
$$
Thus the obstruction asks whether this
residual can be removed by corrections.
It does not ask whether the residual
itself is zero.

For comparison with first-order theory,
put a bar over a polynomial to denote its
class in $A$. The first identity implies
$\overline{F_1}R_0=0$. Consequently the rule
$$
f_i\longmapsto\overline{(F_1)_i}
$$
defines an element of
$\operatorname{Hom}_S(I_0,A)_0$:
it respects every relation among the
$f_i$. This is the embedded tangent space
from
Proposition~%
\ref{proposition-affine-first-order}.

\medskip\noindent
\textbf{The quotient defining cotangent
cohomology.}
Some syzygies exist for any list of
polynomials, simply by commutativity. Let
$K\subset\mathrm{Rel}$ be the submodule
generated by the Koszul syzygies
$$
k_{ij}=f_i\epsilon_j-f_j\epsilon_i.
$$
The quotient $\mathrm{Rel}/K$ is an
$A$-module. To check this, take
$s=\sum_j s_j\epsilon_j\in\mathrm{Rel}$.
Then
$$
f_i s=\sum_j s_j
(f_i\epsilon_j-f_j\epsilon_i)\in K,
$$
since $\sum_j s_jf_j=0$. Thus $I_0$
annihilates $\mathrm{Rel}/K$.

Every $S$-linear map $\phi:E\to A$
vanishes on $K$, because the $f_i$ act as
zero on $A$. Restriction therefore gives
$$
\begin{aligned}
\rho:\operatorname{Hom}_S(E,A)_0
&\longrightarrow
\operatorname{Hom}_A(\mathrm{Rel}/K,A)_0,\\
\rho(\phi)(s+K)&=\phi(s).
\end{aligned}
$$
If $\phi$ is represented by a row $H$ of
polynomials, its values on the chosen
syzygy generators are $\overline{H R_0}$.

\begin{definition}
\label{definition-cotangent-two}
The degree-zero part of the
\emph{second cotangent cohomology} of
$A$ over $k$, with coefficients in $A$, is
$$
T^2_{A,0}=\operatorname{coker}(\rho).
$$
Explicitly,
\begin{equation}
\label{equation-cotangent-two-quotient}
T^2_{A,0}=
\frac{
\operatorname{Hom}_A(\mathrm{Rel}/K,A)_0
}{
\rho\bigl(\operatorname{Hom}_S(E,A)_0\bigr)
}.
\end{equation}
\end{definition}

\noindent
The subscript $0$ means that maps preserve
the usual total degree in the $x_i$.
It does not mean constant polynomial
entries. In our project $E=S(-3)^{16}$,
so a degree-zero map $E\to A$ sends each
basis vector to a cubic in $A$.
The superscript $2$ is the cohomological
index; it is unrelated to this grading.

This is the concrete construction in
\cite[Section 0.2]{Fausk}. Her generator
module is called $F$, her polynomial ring
is $R$, and her Koszul submodule is
$\mathrm{Rel}_0$. Her source for the
restriction map agrees with ours because
$K\subset I_0E$.
The construction and its independence
from the presentation are also explained
in \href{https://math.berkeley.edu/%
\string~robin/math274root.pdf}
{Hartshorne's online lectures}, Section 3,
especially Lemmas 3.2 and 3.3.
These references identify the quotient
with cotangent cohomology; we do not need
to construct a cotangent complex here.

The following theorem explains why this
quotient detects flat lifts at every order.
Write $T^2_A=T^2(A/k,A)$ for the same
cokernel without the degree restriction.

\begin{theorem}[Obstruction to a flat lift]
\label{theorem-obstruction-flat-lift}
Let $A$ be a finitely generated
$k$-algebra. Let $B'\to B$ be a surjection
of local Artinian $k$-algebras with
residue field $k$. Write $J=\ker(B'\to B)$
and assume
$$
\mathfrak{m}_{B'}J=0.
$$
Let $A_B$ be a flat, finitely generated
$B$-algebra, with a specified
identification $A_B\otimes_B k\cong A$.
There is a canonical obstruction
$$
\operatorname{ob}(A_B,B')
\in T^2_A\otimes_k J.
$$
It vanishes if and only if there exists
a flat, finitely generated $B'$-algebra
$A_{B'}$ with an identification
$$
A_{B'}\otimes_{B'}B\cong A_B.
$$
Such a lift can be realized by lifting
the generators of the ideal in a chosen
polynomial presentation of $A_B$.
\end{theorem}

\noindent
For a proof, see Hartshorne,
\href{https://math.berkeley.edu/%
\string~robin/math274root.pdf}
{\emph{Lectures on Deformation Theory}}
(2004), Theorem 10.1(a), p.~58
(PDF page 64; proof on pp.~59--60),
and Corollary 10.3, p.~61.

\medskip\noindent
For the rings $B_r$ above, take
$$
B'=B_{r+1},\qquad B=B_r,\qquad
J=(t^{r+1})/(t^{r+2})\cong k.
$$
Thus a lift extends the specified family
by one order while preserving flatness.
For graded deformations with $\deg t=0$,
the calculation takes place in
$T^2_{A,0}$: the maps preserve degree,
so corrections can be chosen in degree
zero. We now construct the obstruction
explicitly at second order.

\medskip\noindent
\textbf{Why the residual belongs to this
particular target.}
For an arbitrary $s\in\mathrm{Rel}$,
choose a first-order lift $s+tu$ satisfying
$$
F_0u=-F_1s.
$$
Such a $u\in E$ exists because the
first-order deformation is flat. Define
$$
\beta(s)=\overline{F_1u}\in A.
$$
There are three small checks.
\begin{enumerate}
\item If $u'$ is another choice, then
$u'-u\in\mathrm{Rel}$. Since $F_1$ sends
every syzygy into $I_0$, we have
$\overline{F_1(u'-u)}=0$. Thus $\beta(s)$
does not depend on the choice of lift.
Adding or multiplying choices of $u$
also proves that $\beta$ is $S$-linear.

\item For $s=k_{ij}$, take
$$
u=(F_1)_i\epsilon_j-(F_1)_j\epsilon_i.
$$
Then $F_0u=-F_1k_{ij}$, while
$$
F_1u=(F_1)_i(F_1)_j
-(F_1)_j(F_1)_i=0.
$$
Hence $\beta(K)=0$; this reflects the fact
that Koszul syzygies always lift.

\item The degrees in the construction
match. Therefore $\beta$ induces a
degree-zero $A$-linear map
$$
\beta:\mathrm{Rel}/K\longrightarrow A.
$$
\end{enumerate}
On the columns of $R_0$, this map is
represented by $\overline{F_1R_1}$.
The construction proves that these values
respect relations among the syzygies.
An arbitrary row with the same number of
entries need not have this property.

\medskip\noindent
\textbf{Which part can the corrections
remove?}
Reducing the second identity in
\ref{equation-syzygy-lifts} modulo $I_0$
removes $F_0R_2$, and leaves
$$
\beta+\rho(\overline{F_2})=0.
$$
Consequently the obstruction class is
$$
\operatorname{ob}_2(F_1)=[\beta]
\in T^2_{A,0}.
$$
If this class is zero, choose $F_2$ with
$\rho(\overline{F_2})=-\beta$. Each entry
of $F_1R_1+F_2R_0$ then belongs to $I_0$.
Expressing its negative in the generators
$f_i$ supplies the columns of $R_2$.
Thus a second-order flat lift exists.
Conversely, any such lift gives this
identity and hence a zero class.

Notice that $F_2$ is an arbitrary correction
row of the appropriate degrees. It need
not itself represent a first-order
deformation. This is why the denominator
uses all of $\operatorname{Hom}_S(E,A)_0$.

We have already checked independence of
$R_1$ and of the chosen syzygy generators.
For completeness, replacing $F_1$ by
$F_1+F_0H$ merely changes its polynomial
representatives. Taking $R_1-HR_0$ as the
new syzygy correction changes $\beta$ by
$$
-\rho(\overline{F_1H}),
$$
so its class is unchanged. Here $H:E\to E$
has degree zero. Changes of the generator
presentation give the canonical
identifications cited above. Infinitesimal
coordinate changes also preserve the class:
they extend to coordinate changes over
$k[t]/(t^3)$, and the obstruction
construction is natural under these
isomorphisms. Thus the class depends on the
intrinsic first-order deformation.

\begin{example}
\label{example-removing-residual}
Take $S=k[x,y,z]$ and $I_0=(x^2,xy)$.
All syzygies are multiples of the column
$r=(y,-x)^{\mathsf{T}}$: cancellation of
$x$ in $ax^2+bxy=0$ gives $ax+by=0$.
Choose
$$
F_0=(x^2,xy),\quad F_1=(0,yz),\quad
R_0=\begin{pmatrix}y\\-x\end{pmatrix},
\quad R_1=\begin{pmatrix}0\\z\end{pmatrix}.
$$
Then $F_0R_1=xyz$ and $F_1R_0=-xyz$.
The first-order equation holds, but
$$
F_1R_1=yz^2\neq0
$$
even in $A=S/(x^2,xy)$. Nevertheless,
$F_2=(-z^2,0)$ gives $F_2R_0=-yz^2$,
so $R_2=0$ completes the second-order lift.
Indeed the product is identically zero:
$$
(x^2-t^2z^2,xy+tyz)
\begin{pmatrix}y\\-x+tz\end{pmatrix}=0.
$$

Here $\mathrm{Rel}=Sr$ and $K=xSr$.
Giving a map $\mathrm{Rel}/K\to A$ amounts
to choosing its value on $r$, annihilated
by $x$. Since
$$
\operatorname{ann}_A(x)=(x,y),\qquad
\rho(h_1,h_2)(r)=yh_1-xh_2,
$$
restriction is surjective, also in degree
zero: every cubic in $(x,y)\subset A$ has
this form with quadratic $h_1,h_2$.
Thus $T^2_{A,0}=0$. The nonzero residual
represents zero in the quotient.
\end{example}

\medskip\noindent
\textbf{Where the quadratic polynomials
come from.}
The intrinsic tangent space is
$$
V=T^1_{A,0}=
\frac{\operatorname{Hom}_S(I_0,A)_0}
{\operatorname{im}
(\operatorname{Der}_k(S,A)_0
\to\operatorname{Hom}_S(I_0,A)_0)}.
$$
The denominator consists of infinitesimal
coordinate changes. Choose a basis
$v_1,\ldots,v_d$ of $V$ and polynomial
representatives for it. For each basis
vector solve the first-order syzygy
equation once. Linearity then gives
$$
F_1(y)=\sum_{a=1}^d y_a U_a,\qquad
R_1(y)=\sum_{a=1}^d y_a V_a.
$$
Here $U_a$ is a row representing $v_a$,
and $V_a$ is its syzygy correction matrix.
The $y_a$ are coefficients of a tangent
vector, not the original variables $x_i$.

The residual is now
$$
F_1(y)R_1(y)=
\sum_{a,b}y_ay_bU_aV_b.
$$
It is quadratic in $y$. Reduction modulo
$I_0$ and passage to the cokernel are
linear operations. Therefore they give
a quadratic map
$$
q_2:V\longrightarrow T^2_{A,0}.
$$
Choose a basis $e_1,\ldots,e_r$ of
$T^2_{A,0}$. In this basis,
\begin{equation}
\label{equation-quadratic-coordinates}
q_2(y)=\sum_{i=1}^r Q_i(y)e_i,
\qquad Q_i=e_i^*\circ q_2.
\end{equation}
The dual vector $e_i^*$ extracts the
$i$th coordinate. This is literally the
meaning of the polynomial $Q_i$: it is
one coordinate of the obstruction class.

Computationally, choose monomial bases in
the relevant graded pieces of $A$. Encode
maps on syzygies by their coefficients,
impose the relations among syzygies and
the condition that $K$ maps to zero, and
then row-reduce the restriction map $\rho$.
Apply the resulting quotient coordinates
to $\overline{F_1(y)R_1(y)}$. This is
linear elimination producing the $Q_i$.
Equivalently one eliminates the unknown
coefficients of $F_2,R_2$ from the
second-order equation. Identifying the
result with $T^2$ requires the target and
quotient just constructed; the cokernel
of an arbitrary matrix of equations is
not automatically cotangent cohomology.

Changing the basis $e_i$ takes invertible
linear combinations of the $Q_i$.
Changing tangent coordinates substitutes
an invertible linear change in $y$.
The obstruction map itself is intrinsic.
There can be fewer than $r$ independent
quadratic coordinates: the assertion that
all $r$ are independent needs checking.

\medskip\noindent
\textbf{The numbers in the Gr\"unbaum
project.}
The computational source here is
\href{https://github.com/danimalabares/%
grunbaum-zero-context-proof/%
tree/%
ee984e9541e9cb9259728e6dc7b3a1bc0fda8a9b}
{the zero-context proof repository},
at the displayed pinned revision.
Its variables $a,\ldots,h$ use a different
vertex labelling from our
Stanley--Reisner chapter. Direct comparison
of the 16 monomials identifies the ideals
under the substitution
$$
(a,b,c,d,e,f,g,h)\longmapsto
(x_3,x_1,x_7,x_4,x_2,x_6,x_5,x_8).
$$

The audit
\href{https://github.com/danimalabares/%
grunbaum-zero-context-proof/blob/%
ee984e9541e9cb9259728e6dc7b3a1bc0fda8a9b/%
enumerate/SYMMETRY_AUDIT.md}
{\texttt{enumerate/SYMMETRY\_AUDIT.md}}
points to the following exact computations.
Only their tangent and obstruction data
are needed here.
\begin{enumerate}
\item The script
\href{https://github.com/danimalabares/%
grunbaum-zero-context-proof/blob/%
ee984e9541e9cb9259728e6dc7b3a1bc0fda8a9b/%
reconstruct/coordinate_split_rank.m2}
{\texttt{coordinate\_split\_rank.m2}}
computes the normal space and the image
of the 64 linear coordinate changes.
The dimensions are
$$
\dim_k\operatorname{Hom}_S(I_0,A)_0=109,
\qquad \dim_k V=109-56=53.
$$
There are 104 standard cubic monomials in
$A$, so the initial coefficient space has
$16\cdot104=1664$ dimensions. Imposing
the syzygies leaves 109; quotienting by
coordinate changes removes 56 more.
Thus 53 counts intrinsic tangent
directions, not all embedded directions.

\item The script
\href{https://github.com/danimalabares/%
grunbaum-zero-context-proof/blob/%
ee984e9541e9cb9259728e6dc7b3a1bc0fda8a9b/%
reconstruct/audit_ct1_ct2_basis.m2}
{\texttt{audit\_ct1\_ct2\_basis.m2}}
and its
\href{https://github.com/danimalabares/%
grunbaum-zero-context-proof/blob/%
ee984e9541e9cb9259728e6dc7b3a1bc0fda8a9b/%
reconstruct/audit_ct1_ct2_basis_output.txt}
{recorded output} report 53 basis columns
for \texttt{CT\^{}1(\{0\},F0)} and 27
for \texttt{CT\^{}2(\{0\},F0)}.
The latter matrix has size $30\times27$:
its rows refer to the syzygy generators,
and its columns represent a basis of
$$
W=T^2_{A,0},\qquad \dim_k W=27.
$$
The interpretation of \texttt{CT\^{}2}
as cotangent cohomology is specified in
the \href{https://macaulay2.com/doc/%
Macaulay2/share/doc/Macaulay2/%
VersalDeformations/html/___C__T.html}
{Macaulay2 package documentation}.

\item The same Macaulay2 script computes
a fresh order-two versal deformation.
It checks that its 27 quadratic
coordinates have rank 27 and span the
same space as the stored quadratic list.
Separately,
\href{https://github.com/danimalabares/%
grunbaum-zero-context-proof/blob/%
ee984e9541e9cb9259728e6dc7b3a1bc0fda8a9b/%
enumerate/audit_symmetry_formal.py}
{\texttt{audit\_symmetry\_formal.py}}
extracts the degree-two parts of
\texttt{reconstruct/}%
\texttt{p1\_quartic\_component.m2}
and checks their rank over $\mathbb{Q}$.
For example, in that file's coordinates
(with $t_i$ renamed $y_i$), the first two
quadrics are
$$
\begin{aligned}
Q_1&=-y_{40}y_{46}+y_{43}y_{49},\\
Q_2&=y_{11}y_{43}-y_1y_{46}.
\end{aligned}
$$
These are coefficients in the obstruction
space, not additional equations in the
original variables $x_1,\ldots,x_8$.
\end{enumerate}

\noindent
\emph{Verification status (8 September
2026).} The monomial relabelling, the
dimensions 109 and 56, and the rank 27 of
the stored quadrics were recomputed with
exact rational arithmetic while preparing
this section. For the tangent calculation,
all pair syzygies of the monomial ideal
were imposed. For the quadratic rank,
each of the 27 polynomials has a monomial
absent from the other 26; these give a
diagonal coefficient minor with nonzero
entries. Macaulay2 was unavailable, so the
independent computation of
$\dim T^2_{A,0}=27$ and the identification
of the stored span with the freshly
computed primary obstruction coordinates
were checked against the cited scripts
and their recorded output, not rerun.
Those two statements remain computational
inputs to this explanation.

The calculations use rational coefficients;
extension from $\mathbb{Q}$ to
$\mathbb{C}$ preserves these dimensions
and ranks. The 27 independent quadrics
give an injective map
$$
W^*\longrightarrow
\operatorname{Sym}^2(V^*),\qquad
\lambda\longmapsto\lambda\circ q_2.
$$
Thus their span is identified with $W^*$.
Independence is an extra computational
fact, not a consequence of $\dim W=27$.
It also does not say that the quadrics
form a regular sequence or impose
codimension 27.

\medskip\noindent
\textbf{What vanishing proves.}
For a specified tangent vector $y$,
$$
Q_1(y)=\cdots=Q_{27}(y)=0
$$
is equivalent to the existence of a flat
lift from $k[t]/(t^2)$ to $k[t]/(t^3)$.
It does not by itself give an all-orders
deformation or a smooth fibre.

At the next stage, the known residual is
$F_1R_2+F_2R_1$, and the new corrections
are $F_0R_3+F_3R_0$. More generally, after
a flat lift through order $n$, the
coefficient to be removed at order $n+1$
is
$$
\sum_{i=1}^n F_iR_{n+1-i}.
$$
Theorem~\ref{theorem-obstruction-flat-lift}
assigns
its class to the same space $T^2_{A,0}$,
tensored with the one-dimensional kernel
$(t^{n+1})/(t^{n+2})$. These later classes
depend on the previously chosen lift.
The original quadrics describe only the
first of these lifting tests.

\section{Gröbner denerations}
\label{section-grobner-degenerations}

\noindent
Here's the roadmap:
\begin{enumerate}
\item Start with a smooth variety
with algebra $R/I$ that we suspect
might degenerate to the singular
variety.

\item Choose a weight vector
$\omega \in \mathbb{R}^{n+1}$
where $R = \mathbb{C}[x_0,…,x_n]$.
Define
\begin{enumerate}
\item Initial form
$\text{in}_\omega(f) = $
sum of monomials in $f= \sum c_a x^a$
such that $a\cdot \omega$
is maximal.

\item Initial ideal generated
by initial forms $f \in I$.

\item Homoegenization form
$\tilde{f} = \sum c_a x^at^{b-\omega\cdot a}$
where $b = \text{max}\{\omega \cdot a\}$.

\item Homogenization ideal $\tilde{I}$
generated by $\tilde{f}$ for $f \in I$.
\end{enumerate}
\item Apply theorem: then
$R[t]/\tilde{I}$
is flat over $\mathbb{C}[t]$
with fiber $R/\text{in}_\omega(I)$
over $(0)$ and fiber $R/I$ over 
any $(t-u)$ for $u\neq 0$.
\end{enumerate}

\begin{theorem}
\label{theorem-grobner-flat-morphism}
For any ideal $I \subset k[x_1,\ldots,x_r]$,
the $k[t]$-algebra $S[t]/\tilde{I}$ is free
---and thus flat--- as a $k[t]$ module.

[…].

Thus $S[t]/\tilde{I}$ is a flat family over
$k[t]$ of quotients of $S$ whose fiber over
$0$ is $S/\text{in}_\lambda(I)$ and whose
fiber over any $(t-u)$ for $0 \neq u \in k$
is $S/I$.

\end{theorem}

\medskip\noindent
Now I'll put the original notes on this:


\section{Miscellaneous}
\label{section-misc}

Look for the ``Perfect obstruction theory''.
This is used by Bruzzo for deformations of
(super)stacks,
work with E. Paiva and ---.

\bibliography{my}
\bibliographystyle{amsalpha}

\end{document}
